% Source: https://artofproblemsolving.com/community/c3h2154468_how_to_make_any_integer_1_through_22_using_exactly_5_twos_and_some_operations
% Stable identity: aops-topic-2154468-q01
In a puzzle, each integer is to be written using exactly five $2$'s and operations such as $+,-,\times,\div$, powers, square roots, factorials, logarithms, $\lfloor\ \rfloor$ (floor) and $\lceil\ \rceil$ (ceiling). Writing $22$ uses two $2$'s, $222$ uses three, and $.2=0.2$ uses one. A student claims the following. For each claim, find the value of the right-hand side and decide whether the claim is correct (right value and exactly five $2$'s).
\begin{enumerate}
\item[(a)] $29=(2+2)!+2+2+\left\lfloor\sqrt2\right\rfloor$
\item[(b)] $55=\left(2+\dfrac{\left\lfloor\sqrt2\right\rfloor}{.2}\right)\cdot22$
\item[(c)] $105=\dfrac{22+\dfrac22}{.2}$
\item[(d)] $139=\left\lfloor\sqrt{\dfrac{(2\cdot2\cdot2)!}{2}}\right\rfloor-2$
\item[(e)] $152=22\left\lfloor\dfrac{\sqrt2}{2}\right\rfloor-2$
\item[(f)] $197=\left\lceil22\log_{\sqrt2}22\right\rceil$
\item[(g)] $215=222+\left\lfloor\dfrac{\sqrt2}{.2}\right\rfloor$
\end{enumerate}
Then correct (c) and (g) by changing one symbol in each.
